Diffusion models \citep{pmlr-v37-sohl-dickstein15} are latent-variable generative models of the form
\(
p_\theta(\mathbf{x}_0)
=
\int p_\theta(\mathbf{x}_{0:T})\, d\mathbf{x}_{1:T},
\)
where \(\mathbf{x}_{1:T}\) are latents of the same dimensionality as the observed data
\(\mathbf{x}_0\). Following \citet{NEURIPS2020_4c5bcfec}, the model consists of a fixed
\emph{forward diffusion process} that gradually adds noise to data, and a learned
\emph{reverse process} that aims to invert this corruption.

\begin{figure}[H]
    \centering
    \includegraphics[width=0.85\textwidth]{chapters/background/diffusion_models/media/DDPM.png}
    \caption{Directed graphical model structure of DDPMs \citep{NEURIPS2020_4c5bcfec}.}
    \label{fig:ddpm_graphical_model}
\end{figure}

The joint distribution of all variables is written as a Markov chain:
\begin{equation}
p_\theta(\mathbf{x}_{0:T})
=
p(\mathbf{x}_T)
\prod_{t=1}^{T} p_\theta(\mathbf{x}_{t-1}\mid \mathbf{x}_t),
\qquad
p(\mathbf{x}_T)=\mathcal{N}(\mathbf{0},\mathbf{I}),
\label{eq:ddpm_joint}
\end{equation}
where \(p_\theta(\mathbf{x}_{t-1}\mid\mathbf{x}_t)\) is the reverse process. The forward diffusion
process is defined as a separate Markov chain:
\begin{equation}
q(\mathbf{x}_{1:T}\mid\mathbf{x}_0)
=
\prod_{t=1}^T
q(\mathbf{x}_t\mid\mathbf{x}_{t-1}),
\qquad
q(\mathbf{x}_t\mid\mathbf{x}_{t-1})
=
\mathcal{N}\!\left(
\sqrt{1-\beta_t}\,\mathbf{x}_{t-1},\;
\beta_t \mathbf{I}
\right),
\label{eq:ddpm_forward}
\end{equation}
where \(\{\beta_t\}_{t=1}^T\) is the variance schedule that \emph{can} be learned by reparameterization,
or fixed to predetermined constants (e.g., a linear schedule from $10^{-4}$ to $10^{-2}$) such that \(\beta_t\) typically increases slowly with \(t\). 
This ensures that the forward chain remains
analytically tractable and preserves the same functional form across timesteps while gradually
driving $\mathbf{x}_T$ toward an isotropic Gaussian.

The authors \citep{NEURIPS2020_4c5bcfec} choose \(T\) large (e.g.\ \(T=1000\)) and let \(\beta_t\) increase linearly or according
to a smooth curve from a very small value to a modest value. This ensures that each step only
slightly degrades the signal, so the forward chain remains close to a diffusion process with
well-behaved Gaussian conditionals, and by the final step \(\mathbf{x}_T\) is nearly standard
Gaussian.

%%%%%%%%%%%%%%%%%%%%%%%%%%%%%%%%%%%%%%%%%%%%%%%%%%%%%%%%%%%%%%%%%%%%%%%%%%%%%%%
\subparagraph{Forward Diffusion Process.}
%%%%%%%%%%%%%%%%%%%%%%%%%%%%%%%%%%%%%%%%%%%%%%%%%%%%%%%%%%%%%%%%%%%%%%%%%%%%%%%

Because each transition in \eqref{eq:ddpm_forward} is linear Gaussian, the forward chain admits a
closed-form marginal at any timestep and repeatedly composing the Gaussian transitions yields:
\begin{equation}
q(\mathbf{x}_t\mid\mathbf{x}_0)
=
\mathcal{N}\!\left(
\sqrt{\bar\alpha_t}\,\mathbf{x}_0,\;
(1-\bar\alpha_t)\mathbf{I}
\right) 
\text{ where} \quad
\alpha_t := 1-\beta_t,
\qquad
\bar\alpha_t := \prod_{s=1}^{t}\alpha_s.
\label{eq:ddpm_forward_marginal}
\end{equation}
We can define $x_t$ through an explicit reparameterization
\begin{align}
\mathbf{x}_t &= \sqrt{\alpha_t} \, \mathbf{x}_{t-1}
+
\sqrt{1-\alpha_t}\,\boldsymbol{\epsilon}_{t-2} \\
&= \sqrt{\alpha\,\alpha_{t-1}}\, \mathbf{x}_{t-2} 
+
\sqrt{1-\alpha_t\,1-\alpha_{t-1}}\,\boldsymbol{\epsilon}_{t-1} \\
&= \dots \\
&=
\sqrt{\bar\alpha_t}\,\mathbf{x}_0
+
\sqrt{1-\bar\alpha_t}\,\boldsymbol{\epsilon},
\qquad
\boldsymbol{\epsilon}\sim\mathcal{N}(\mathbf{0},\mathbf{I}).
\label{eq:reparameterization_of_xt}
\end{align}
where the state at timestep \(t\) is a combination of the original data and Gaussian noise,
with the signal coefficient \(\sqrt{\bar\alpha_t}\) gradually shrinking as \(t\) increases.



%%%%%%%%%%%%%%%%%%%%%%%%%%%%%%%%%%%%%%%%%%%%%%%%%%%%%%%%%%%%%%%%%%%%%%%%%%%%%%%
\subparagraph{Reverse Denoising Process.}
%%%%%%%%%%%%%%%%%%%%%%%%%%%%%%%%%%%%%%%%%%%%%%%%%%%%%%%%%%%%%%%%%%%%%%%%%%%%%%%

The reverse process attempts to invert the forward diffusion by learning Gaussian transitions of the
form
\begin{equation}
p_\theta(\mathbf{x}_{t-1}\mid\mathbf{x}_t)
=
\mathcal{N}\!\left(
\boldsymbol{\mu}_\theta(\mathbf{x}_t,t),\;
\boldsymbol{\Sigma}_\theta(\mathbf{x}_t,t)
\right).
\label{eq:ddpm_reverse}
\end{equation}
Here, \(\theta\) denotes a single shared neural-network parameterization; the dependence on \(t\) is
provided through a timestep embedding. Because the forward process $q(\mathbf{x}_{1:T}\mid\mathbf{x}_0)$ is fixed and Gaussian, the
associated conditional distributions $q(\mathbf{x}_{t-1}\mid\mathbf{x}_t,\mathbf{x}_0)$ admit
closed-form expressions. These conditionals arise directly from standard Gaussian identities:
they describe how the forward noising process relates adjacent latent variables when conditioned
on the original clean sample $\mathbf{x}_0$. The DDPM model parameterizes
$p_\theta(\mathbf{x}_{t-1}\mid\mathbf{x}_t)$ to approximate these analytically derived
Gaussian conditionals.
It has the Gaussian form
\begin{equation}
q(\mathbf{x}_{t-1}\mid\mathbf{x}_t,\mathbf{x}_0)
=
\mathcal{N}\!\left(
\tilde{\boldsymbol{\mu}}_t(\mathbf{x}_t,\mathbf{x}_0),\;
\tilde{\beta}_t \mathbf{I}
\right),
\label{eq:true_reverse_posterior}
\end{equation}
where
\begin{align}
\tilde{\boldsymbol{\mu}}_t(\mathbf{x}_t,\mathbf{x}_0)
&=
\frac{\sqrt{\bar\alpha_{t-1}}\beta_t}{1-\bar\alpha_t}\,\mathbf{x}_0
+
\frac{\sqrt{\alpha_t}(1-\bar\alpha_{t-1})}{1-\bar\alpha_t}\,\mathbf{x}_t,
\\[4pt]
\tilde{\beta}_t
&=
\frac{1-\bar\alpha_{t-1}}{1-\bar\alpha_t}\,\beta_t.
\end{align}
The model therefore learns $p_\theta(\mathbf{x}_{t-1}\mid\mathbf{x}_t)$ so that its mean matches
the analytically derived forward-process posterior. For stability, DDPM fixes the reverse-process
variance to one of two experimentally evaluated choices:
\[
\sigma_t^2 = \beta_t
\qquad\text{or}\qquad
\sigma_t^2 = \tilde{\beta}_t = \frac{1-\bar{\alpha}_{t-1}}{1-\bar{\alpha}_t}\,\beta_t,
\]
as reported by \citet{NEURIPS2020_4c5bcfec}. Both forms preserve the Gaussian functional structure
needed for closed-form KL divergences in the ELBO, while yielding similar empirical performance.


%%%%%%%%%%%%%%%%%%%%%%%%%%%%%%%%%%%%%%%%%%%%%%%%%%%%%%%%%%%%%%%%%%%%%%%%%%%%%%%
\subparagraph{Training Objective.}
%%%%%%%%%%%%%%%%%%%%%%%%%%%%%%%%%%%%%%%%%%%%%%%%%%%%%%%%%%%%%%%%%%%%%%%%%%%%%%%

Training minimizes the negative log-likelihood \(-\log p_\theta(\mathbf{x}_0)\). Since
\(\mathbf{x}_{1:T}\) is intractable to marginalize exactly, one optimizes a standard variational
lower bound:
\begin{equation}
\mathbb{E}_q[-\log p_\theta(\mathbf{x}_0)]
\;\le\;
\mathcal{L}(\theta)
=
\mathbb{E}_q\!\left[
-\log\frac{p_\theta(\mathbf{x}_{0:T})}{q(\mathbf{x}_{1:T}\mid\mathbf{x}_0)}
\right]
=
\mathbb{E}_q\!\left[
-\log p(\mathbf{x}_T)
-
\sum_{t=1}^T
\log\frac{p_\theta(\mathbf{x}_{t-1}\mid\mathbf{x}_t)}
{q(\mathbf{x}_t\mid\mathbf{x}_{t-1},\mathbf{x}_0)}
\right]
=:
L(\theta).
\label{eq:elbo_ddpm_raw}
\end{equation}
This expression can be rewritten as a sum of KL terms \citep{weng2021diffusion}:
\begin{equation}
\mathcal{L}(\theta)
=
\mathbb{E}_q\Big[
\underbrace{
D_{\mathrm{KL}}\bigl(q(\mathbf{x}_T\mid\mathbf{x}_0)\,\|\,p(\mathbf{x}_T)\bigr)
}_{L_T}
+
\sum_{t=2}^T
\underbrace{
D_{\mathrm{KL}}\bigl(
q(\mathbf{x}_{t-1}\mid\mathbf{x}_t,\mathbf{x}_0)
\;\|\;
p_\theta(\mathbf{x}_{t-1}\mid\mathbf{x}_t)
\bigr)
}_{L_{t-1}}
-
\underbrace{\log p_\theta(\mathbf{x}_0\mid\mathbf{x}_1)}_{L_0}
\Big].
\label{eq:elbo_ddpm_kl}
\end{equation}
The term \(L_T\) encourages the forward terminal distribution to match the prior. Each \(L_{t-1}\)
penalizes discrepancies between the true reverse posterior and the learned reverse conditional at
that timestep. The final term \(L_0\) plays the role of a reconstruction likelihood for
\(\mathbf{x}_0\). Because all of these terms compare Gaussian distributions with known parameters,
they admit closed-form expressions, allowing efficient estimation with low-variance gradients.

%%%%%%%%%%%%%%%%%%%%%%%%%%%%%%%%%%%%%%%%%%%%%%%%%%%%%%%%%%%%%%%%%%%%%%%%%%%%%%%
%%%%%%%%%%%%%%%%%%%%%%%%%%%%%%%%%%%%%%%%%%%%%%%%%%%%%%%%%%%%%%%%%%%%%%%%%%%%%%%
\subparagraph{Simplified Noise-Prediction Objective.}
%%%%%%%%%%%%%%%%%%%%%%%%%%%%%%%%%%%%%%%%%%%%%%%%%%%%%%%%%%%%%%%%%%%%%%%%%%%%%%%

The variational bound in \eqref{eq:elbo_ddpm_kl} has the generic form
\(
L_{\mathrm{VLB}} = L_T + \sum_{t=2}^T L_{t-1} + L_0,
\)
where $L_T$ and $L_0$ do not depend on the parametrization of the reverse mean.
The terms $L_{t-1}$ are KL divergences between Gaussians and therefore can be written in closed
form. Assuming an isotropic reverse covariance
$\boldsymbol{\Sigma}_\theta(\mathbf{x}_t,t)=\sigma_t^2\mathbf{I}$, each such term reduces to
a mean-squared error between the true posterior mean $\tilde{\boldsymbol{\mu}}_t$ and the model
mean $\boldsymbol{\mu}_\theta$:
\begin{equation}
L_{t-1}
=
\mathbb{E}_q\!\left[
\frac{1}{2\sigma_t^2}
\bigl\|
\tilde{\boldsymbol{\mu}}_t(\mathbf{x}_t,\mathbf{x}_0)
-
\boldsymbol{\mu}_\theta(\mathbf{x}_t,t)
\bigr\|^2
\right]
+ C_t,
\label{eq:L_t_mse_form}
\end{equation}
where $C_t$ does not depend on $\theta$.

Using the closed-form forward posterior
\(
q(\mathbf{x}_{t-1}\mid\mathbf{x}_t,\mathbf{x}_0)
=
\mathcal{N}(\tilde{\boldsymbol{\mu}}_t,\tilde{\beta}_t\mathbf{I})
\)
and the forward reparameterization
\(
\mathbf{x}_t = \sqrt{\bar{\alpha}_t}\,\mathbf{x}_0 +
\sqrt{1-\bar{\alpha}_t}\,\boldsymbol{\epsilon},\;
\boldsymbol{\epsilon}\sim\mathcal{N}(0,I),
\)
the posterior mean can be expressed as
\begin{equation}
\tilde{\boldsymbol{\mu}}_t(\mathbf{x}_t,\mathbf{x}_0)
=
\frac{1}{\sqrt{\alpha_t}}
\left(
\mathbf{x}_t
-
\frac{1-\alpha_t}{\sqrt{1-\bar{\alpha}_t}}\,
\boldsymbol{\epsilon}
\right),
\end{equation}
Since $\mathbf{x}_t$ is available as model
input, it is convenient to parametrize the reverse mean so that the network predicts the noise
directly:
\begin{equation}
\boldsymbol{\mu}_\theta(\mathbf{x}_t,t)
=
\frac{1}{\sqrt{\alpha_t}}
\left(
\mathbf{x}_t
-
\frac{1-\alpha_t}{\sqrt{1-\bar{\alpha}_t}}\,
\boldsymbol{\epsilon}_\theta(\mathbf{x}_t,t)
\right).
\label{eq:reverse_mean_eps_param}
\end{equation}
Substituting $\tilde{\boldsymbol{\mu}}_t$ and $\boldsymbol{\mu}_\theta$ into
\eqref{eq:L_t_mse_form} and simplifying gives
\begin{equation}
L_{t-1}
=
\mathbb{E}_{\mathbf{x}_0,\boldsymbol{\epsilon}}
\left[
\frac{(1-\alpha_t)^2}{2\sigma_t^2\,\alpha_t(1-\bar{\alpha}_t)}
\,
\bigl\|
\boldsymbol{\epsilon}
-
\boldsymbol{\epsilon}_\theta(\mathbf{x}_t,t)
\bigr\|^2
\right]
+ C_t,
\end{equation}
so each KL term is a weighted mean-squared error between the true noise
$\boldsymbol{\epsilon}$ and the predicted noise $\boldsymbol{\epsilon}_\theta(\mathbf{x}_t,t)$.

Dropping the timestep-dependent weight yields the commonly used simplified objective
\begin{equation}
L_{\mathrm{simple}}(\theta)
=
\mathbb{E}_{t,\mathbf{x}_0,\boldsymbol{\epsilon}}
\bigl[
\|
\boldsymbol{\epsilon}
-
\boldsymbol{\epsilon}_\theta(\mathbf{x}_t,t)
\|^2
\bigr],
\qquad
\mathbf{x}_t
=
\sqrt{\bar{\alpha}_t}\,\mathbf{x}_0
+
\sqrt{1-\bar{\alpha}_t}\,\boldsymbol{\epsilon},
\label{eq:ddpm_simple_loss}
\end{equation}
which is equivalent to a reweighted form of the ELBO but is simpler to implement and has been
observed to yield improved training stability and sample quality in practice.


%%%%%%%%%%%%%%%%%%%%%%%%%%%%%%%%%%%%%%%%%%%%%%%%%%%%%%%%%%%%%%%%%%%%%%%%%%%%%%%
\subparagraph{Training and Sampling Algorithms.}

The training and sampling procedures are summarized in Figure~\ref{fig:algorithms}. Training draws
a data sample \(\mathbf{x}_0\), chooses a timestep \(t\) uniformly from \(\{1,\dots,T\}\), samples
\(\boldsymbol{\epsilon}\sim\mathcal{N}(0,I)\), constructs \(\mathbf{x}_t\) using the forward
marginal \eqref{eq:ddpm_forward_marginal}, and takes a gradient step on the objective
\eqref{eq:ddpm_simple_loss}. Sampling inverts this process: starting from
\(\mathbf{x}_T\sim\mathcal{N}(\mathbf{0},\mathbf{I})\), the model repeatedly applies the reverse
transition implied by \eqref{eq:reverse_mean_eps_param},
\[
\mathbf{x}_{t-1}
=
\frac{1}{\sqrt{\alpha_t}}
\left(
\mathbf{x}_t
-
\frac{\beta_t}{\sqrt{1-\bar\alpha_t}}\,
\boldsymbol{\epsilon}_\theta(\mathbf{x}_t,t)
\right)
+
\sigma_t \mathbf{z},
\qquad
\mathbf{z}\sim\mathcal{N}(\mathbf{0},\mathbf{I}),
\]
for \(t = T,T-1,\dots,1\), where \(\sigma_t^2\) is chosen consistently with the forward variance
schedule. This iterative denoising procedure gradually transforms white noise into a structured
sample \(\mathbf{x}_0\).

\begin{figure}[H]
    \centering
    \includegraphics[width=0.85\textwidth]{chapters/background/diffusion_models/media/Algorithms.png}
    \caption{
        \textbf{DDPM Training and Sampling Algorithms.}
        Algorithm~1 implements stochastic noise-prediction training based on the closed-form forward
        marginal: at each iteration, a random timestep and noise level are chosen, and the network
        learns to predict the noise that produced the corresponding corrupted sample. Algorithm~2
        applies the learned reverse transitions from a Gaussian prior $\mathbf{x}_T$ down to
        $\mathbf{x}_0$, yielding new samples by successive denoising \citep{NEURIPS2020_4c5bcfec}.
    }
    \label{fig:algorithms}
\end{figure}

%%%%%%%%%%%%%%%%%%%%%%%%%%%%%%%%%%%%%%%%%%%%%%%%%%%%%%%%%%%%%%%%%%%%%%%%%%%%%%%
\subparagraph{Architecture.}

The reverse-process mean or noise predictor is implemented using a U-Net architecture with residual
blocks, skip connections, and multi-scale feature hierarchies \citep{NEURIPS2020_4c5bcfec}. The
timestep \(t\) is encoded by a sinusoidal positional embedding and injected into intermediate layers,
allowing the network to adapt its predictions to the current noise level. Coarse scales capture
global semantic structure, while finer scales refine local details. The network outputs
\(\boldsymbol{\epsilon}_\theta(\mathbf{x}_t,t)\), which is converted to the reverse mean via
\eqref{eq:reverse_mean_eps_param}, and the variance is typically fixed as
\(\boldsymbol{\Sigma}_\theta(\mathbf{x}_t,t)=\sigma_t^2\mathbf{I}\) to keep training stable.

%%%%%%%%%%%%%%%%%%%%%%%%%%%%%%%%%%%%%%%%%%%%%%%%%%%%%%%%%%%%%%%%%%%%%%%%%%%%%%%
% % \subparagraph{Summary.}

% % Denoising Diffusion Probabilistic Models treat generative modeling as the inversion of a fixed
% % Gaussian diffusion process. A carefully chosen small variance schedule ensures that the forward
% % chain corrupts the data gradually while leading to a nearly Gaussian terminal distribution. The
% % reverse chain is learned via a variational objective whose dominant contribution can be expressed as
% % a simple noise-prediction loss. The resulting sampling procedure recovers data by iterative
% denoising, forming the foundation for many modern diffusion-based generative models.
